\section{How General the Switch Is}\label{sec:general}

\subsection{Other Benchmarks}\label{sec:ood}

We ask whether removing the shared region from adapters trained only on MATH also turns thinking off on other problems. We test this on GSM8K \citep{cobbe2021training}, an easier set of math word problems, and on GPQA-Diamond \citep{rein2023gpqa}, a multiple-choice science benchmark outside math.

Removing the shared region lowers reflective markers to 0.2--9\% of their original level (9--39 per 1,000 words) on GSM8K and to 0.1--10\% of their original level (9--31) on GPQA-Diamond, whereas removing the same energy from other directions outside the shared region leaves 50--197\% and 27--140\%, respectively (Table~\ref{tab:12}).
On GSM8K, responses become 5--19 times shorter, and the share of responses that hit the generation limit drops from 11--84\% to 2\% or less. Accuracy rises from a nine-cell mean of 34.2\% to 47.6\%, close to the $\lambda=1$ adapter trained on short solutions (48.5\%). Removing the same energy from other directions reaches only 39.3\%, lower than removing the shared region in all nine cells (Table~\ref{tab:13}). B/Llama is the only cell whose accuracy stays the same when the shared region is removed.
On GPQA-Diamond, accuracy under every condition is below chance (25\%): the means are 13.7--20.1\%, and even the $\lambda=1$ adapter trained on short solutions reaches only 22.0\%. We therefore read only reflective markers and length on this benchmark (Table~\ref{tab:12}). For OLMo, whose context is 4,096 tokens, some questions are long enough that the generation limit shrinks to 1,293 tokens, so its responses become only 1.5--1.6 times shorter (3.8--4.9 times for the other two students).

Because a region found only from MATH problems turns thinking off in the same way on math problems of a different difficulty and format and on science problems outside math, the region appears to be tied to the way of solving by thinking rather than to the content of MATH problems. The only evidence from outside math, however, is GPQA-Diamond, and even there we judge by reflective markers and length rather than accuracy. In addition, we did not run the length-matched cut outside MATH, so we do not separate how much of the GSM8K accuracy gain comes from resolving truncation.
In short, removing the shared region turns thinking off and lets the student finish short solutions on math and science problems outside MATH as well, so the switch works beyond the problems it was trained on.

\begin{table*}[t]
\centering
\small
% Table 12 (wide: use table*), generated from overleaf/06_how_general_the_switch_is.md
\resizebox{\linewidth}{!}{%
\begin{tabular}{llrrrrrr}
\toprule
Pair & Student & GSM8K: $\lambda$ = 0 & GSM8K: shared removed & GSM8K: other removed & GPQA-Diamond: $\lambda$ = 0 & GPQA-Diamond: shared removed & GPQA-Diamond: other removed \\
\midrule
A & Qwen2.5-0.5B & 28.7 $\pm$ 4.0 & 0.41 $\pm$ 0.18 & 24.4 $\pm$ 7.1 & 13.3 $\pm$ 2.1 & 0.16 $\pm$ 0.15 & 12.3 $\pm$ 2.9 \\
 & Llama-3.2-1B & 39.1 $\pm$ 1.8 & 1.0 $\pm$ 1.0 & 19.5 $\pm$ 0.89 & 31.5 $\pm$ 1.2 & 0.04 $\pm$ 0.03 & 17.2 $\pm$ 1.4 \\
 & OLMo-2-1B & 20.1 $\pm$ 1.3 & 1.0 $\pm$ 0.41 & 12.4 $\pm$ 2.0 & 9.1 $\pm$ 0.64 & 0.72 $\pm$ 0.03 & 5.1 $\pm$ 1.1 \\
\midrule
B & Qwen2.5-0.5B & 9.0 $\pm$ 1.3 & 0.29 $\pm$ 0.23 & 17.6 $\pm$ 7.3 & 10.6 $\pm$ 0.88 & 0.06 $\pm$ 0.01 & 6.5 $\pm$ 1.6 \\
 & Llama-3.2-1B & 15.9 $\pm$ 8.0 & 0.09 $\pm$ 0.15 & 13.0 $\pm$ 3.8 & 25.0 $\pm$ 7.9 & 0.04 $\pm$ 0.02 & 6.8 $\pm$ 1.9 \\
 & OLMo-2-1B & 10.2 $\pm$ 2.8 & 0.94 $\pm$ 0.13 & 8.6 $\pm$ 1.5 & 9.2 $\pm$ 0.76 & 0.95 $\pm$ 0.09 & 5.8 $\pm$ 1.2 \\
\midrule
C & Qwen2.5-0.5B & 20.9 $\pm$ 3.7 & 0.19 $\pm$ 0.13 & 28.5 $\pm$ 3.1 & 9.0 $\pm$ 1.5 & 0.03 $\pm$ 0.04 & 12.7 $\pm$ 2.9 \\
 & Llama-3.2-1B & 22.1 $\pm$ 2.7 & 0.05 $\pm$ 0.01 & 23.9 $\pm$ 2.9 & 22.0 $\pm$ 0.38 & 0.02 $\pm$ 0.01 & 25.3 $\pm$ 1.1 \\
 & OLMo-2-1B & 21.8 $\pm$ 2.8 & 0.61 $\pm$ 0.02 & 15.7 $\pm$ 1.8 & 9.2 $\pm$ 0.46 & 0.83 $\pm$ 0.07 & 7.8 $\pm$ 0.90 \\
\midrule
Mean tokens &  & 3,596 & 353 & 3,268 & 3,319 & 938 & 2,687 \\
\midrule
Mean acc. (\%) &  & 34.2 & 47.6 & 39.3 & 13.7 & 20.1 & 17.6 \\
\bottomrule
\end{tabular}}

\caption{Frequency of reflective markers (per 1,000 words) on GSM8K and GPQA-Diamond when the shared region is removed from the $\lambda=0$ adapter trained on MATH; the last two rows give the nine-cell mean token count and accuracy (\%). Other removed removes the same energy from other directions outside the shared region. Outside MATH as well, reflective markers vanish only when the shared region is removed. Values after $\pm$ are standard deviations across three seeds.}
\label{tab:12}
\end{table*}

\begin{table*}[t]
\centering
\small
% Table 13 (fits one column or use table*), generated from overleaf/06_how_general_the_switch_is.md
\begin{tabular}{llrrrrr}
\toprule
Pair & Student & $\lambda$ = 0 & Shared removed & Other removed & $\lambda$ = 1 & Truncation rate ($\lambda$ = 0 $\rightarrow$ shared removed) \\
\midrule
A & Qwen2.5-0.5B & 25.3 $\pm$ 1.8 & 33.4 $\pm$ 5.5 & 31.6 $\pm$ 2.5 & 48.3 $\pm$ 2.3 & 0.60 $\pm$ 0.01 $\rightarrow$ 0.02 $\pm$ 0.01 \\
 & Llama-3.2-1B & 17.7 $\pm$ 2.6 & 34.1 $\pm$ 2.0 & 28.8 $\pm$ 1.8 & 40.7 $\pm$ 1.5 & 0.84 $\pm$ 0.02 $\rightarrow$ 0.01 $\pm$ 0.01 \\
 & OLMo-2-1B & 43.8 $\pm$ 2.0 & 66.6 $\pm$ 0.4 & 58.1 $\pm$ 2.3 & 60.8 $\pm$ 1.3 & 0.42 $\pm$ 0.03 $\rightarrow$ 0.01 $\pm$ 0.00 \\
\midrule
B & Qwen2.5-0.5B & 34.2 $\pm$ 2.3 & 44.3 $\pm$ 0.5 & 37.2 $\pm$ 3.5 & 43.2 $\pm$ 0.6 & 0.21 $\pm$ 0.02 $\rightarrow$ 0.01 $\pm$ 0.00 \\
 & Llama-3.2-1B & 36.5 $\pm$ 1.3 & 36.5 $\pm$ 1.5 & 36.2 $\pm$ 3.1 & 42.8 $\pm$ 1.0 & 0.18 $\pm$ 0.01 $\rightarrow$ 0.01 $\pm$ 0.00 \\
 & OLMo-2-1B & 55.6 $\pm$ 2.0 & 67.7 $\pm$ 1.4 & 58.8 $\pm$ 0.9 & 57.7 $\pm$ 0.9 & 0.11 $\pm$ 0.01 $\rightarrow$ 0.01 $\pm$ 0.01 \\
\midrule
C & Qwen2.5-0.5B & 27.9 $\pm$ 0.6 & 45.3 $\pm$ 2.7 & 31.9 $\pm$ 2.7 & 39.8 $\pm$ 1.4 & 0.46 $\pm$ 0.01 $\rightarrow$ 0.01 $\pm$ 0.00 \\
 & Llama-3.2-1B & 26.0 $\pm$ 4.9 & 34.6 $\pm$ 0.7 & 21.2 $\pm$ 1.4 & 44.1 $\pm$ 0.3 & 0.64 $\pm$ 0.03 $\rightarrow$ 0.00 $\pm$ 0.00 \\
 & OLMo-2-1B & 40.9 $\pm$ 5.0 & 65.9 $\pm$ 1.2 & 49.8 $\pm$ 3.9 & 59.3 $\pm$ 0.5 & 0.51 $\pm$ 0.03 $\rightarrow$ 0.01 $\pm$ 0.01 \\
\midrule
Mean &  & 34.2 & 47.6 & 39.3 & 48.5 & 0.44 $\rightarrow$ 0.01 \\
\bottomrule
\end{tabular}

\caption{GSM8K accuracy (\%) under the same conditions, and the share of responses that hit the generation limit. Removing the shared region almost eliminates truncation and raises accuracy to the level of the $\lambda=1$ adapter trained on short solutions. Values after $\pm$ are standard deviations across three seeds.}
\label{tab:13}
\end{table*}

\subsection{Without Interpolated Teachers}\label{sec:noninterp}

We next ask whether the shared region is an artifact of interpolating teachers or also arises in ordinary reasoning distillation without interpolation. We therefore find the region from adapters trained with the same reasoning teacher ($\lambda=0$) that differ only in seed (the two-seed region) and remove it. The adapter being edited is not used to build the region, which avoids the circularity of removing a region built from the adapter itself.

The two-seed region covers 54--60\% of the dimensions of the shared region. Removing it lowers reflective markers from a nine-cell mean of 22 to 1.7 per 1,000 words and makes responses 3.4 times shorter, whereas removing the same energy from other directions outside this region leaves 18.2, and responses keep 79--113\% of their original length (Table~\ref{tab:14}).
Other methods that turn thinking off without additional training leave more reflective markers. Removing the top singular components of the update by the criterion of Gold-Switch ($\tau=0.6$) \citep{lee2025goldswitch} leaves 5.2, inserting a sentence stating that thinking is finished \citep{ma2025nothinking} leaves 16.2, and a length-matched cut \citep{muennighoff2025s1} leaves 15.0. Merging half-and-half with the short-solution adapter \citep{wu2025merging} leaves 3.5, but it requires a separate adapter trained on short-solution data (Table~\ref{tab:14}).
Accuracy rises from a nine-cell mean of 13.9\% to 21.7\%, the highest among the methods without additional training, but this only brings it back to a level similar to or below the untrained student (Qwen 35.1\%, Llama 23.3\%, OLMo 22.2\%) (Table~\ref{tab:14}).

Because the region that turns thinking on and off is captured to a large extent by just two adapters trained without interpolated teachers, it is not merely a product of the interpolation design but also arises in reasoning distillation without interpolation. Two seeds, however, capture only a little over half of the shared region, and removing it is weaker than removing the original region (0.07--1.5 reflective markers per 1,000 words, 24.7\% accuracy).
That the $\lambda=0$ adapter itself falls below the untrained student is largely because 51--87\% of its responses hit the generation limit, which is in line with reports that small students perform worse when they learn long solutions \citep{li2025small,luo2025valley}. Since the model with thinking turned off is not better than the untrained student either, this result should not be read as a way to build a better model.
In short, a little over half of this region is captured by just two reasoning-distillation adapters trained without interpolated teachers, and removing that part alone turns thinking off.

\begin{table*}[t]
\centering
\small
% Table 14 (wide: use table*), generated from overleaf/06_how_general_the_switch_is.md
\resizebox{\linewidth}{!}{%
\begin{tabular}{llrrrrrrrrr}
\toprule
Pair & Student & $\lambda$ = 0 & Seed region removed & Length-matched cut & Other removed (seed region) & Top removed ($\tau$ = 0.6) & Think-skip prompt & Merged with $\lambda$ = 1 & $\lambda$ = 1 & Untrained \\
\midrule
A & Qwen2.5-0.5B & 15.2 $\pm$ 1.4 & 18.2 $\pm$ 2.6 & 18.8 $\pm$ 2.0 & 17.0 $\pm$ 2.0 & 17.1 $\pm$ 1.2 & 15.8 $\pm$ 1.4 & 34.8 $\pm$ 0.7 & 34.4 $\pm$ 1.0 & 35.1 $\pm$ 0.9 \\
 & Llama-3.2-1B & 11.8 $\pm$ 1.3 & 18.8 $\pm$ 1.8 & 21.0 $\pm$ 0.9 & 13.6 $\pm$ 2.9 & 17.1 $\pm$ 2.3 & 10.8 $\pm$ 0.3 & 29.4 $\pm$ 1.2 & 29.8 $\pm$ 3.5 & 23.3 $\pm$ 0.4 \\
 & OLMo-2-1B & 10.2 $\pm$ 2.3 & 23.8 $\pm$ 0.7 & 11.2 $\pm$ 1.4 & 13.8 $\pm$ 2.0 & 20.7 $\pm$ 0.9 & 10.2 $\pm$ 1.8 & 18.4 $\pm$ 1.5 & 19.8 $\pm$ 2.2 & 22.2 $\pm$ 0.4 \\
\midrule
B & Qwen2.5-0.5B & 17.5 $\pm$ 0.7 & 26.7 $\pm$ 6.3 & 14.7 $\pm$ 1.2 & 19.9 $\pm$ 1.9 & 16.4 $\pm$ 2.9 & 16.9 $\pm$ 0.9 & 29.6 $\pm$ 0.9 & 26.2 $\pm$ 1.6 & 35.1 $\pm$ 0.9 \\
 & Llama-3.2-1B & 17.1 $\pm$ 0.9 & 25.1 $\pm$ 0.5 & 20.1 $\pm$ 0.8 & 20.0 $\pm$ 3.4 & 18.8 $\pm$ 1.8 & 19.2 $\pm$ 1.9 & 28.3 $\pm$ 0.4 & 23.6 $\pm$ 1.3 & 23.3 $\pm$ 0.4 \\
 & OLMo-2-1B & 10.5 $\pm$ 2.0 & 22.0 $\pm$ 0.7 & 11.2 $\pm$ 0.1 & 16.8 $\pm$ 2.0 & 20.9 $\pm$ 1.0 & 10.4 $\pm$ 1.3 & 17.7 $\pm$ 1.0 & 16.2 $\pm$ 1.1 & 22.2 $\pm$ 0.4 \\
\midrule
C & Qwen2.5-0.5B & 16.2 $\pm$ 1.1 & 20.2 $\pm$ 1.3 & 17.4 $\pm$ 0.7 & 16.0 $\pm$ 0.8 & 20.2 $\pm$ 1.3 & 16.6 $\pm$ 2.0 & 24.5 $\pm$ 2.0 & 24.9 $\pm$ 0.9 & 35.1 $\pm$ 0.9 \\
 & Llama-3.2-1B & 15.0 $\pm$ 2.0 & 20.5 $\pm$ 0.7 & 19.5 $\pm$ 2.1 & 14.2 $\pm$ 0.4 & 18.9 $\pm$ 2.3 & 14.1 $\pm$ 3.9 & 31.2 $\pm$ 0.8 & 30.3 $\pm$ 1.7 & 23.3 $\pm$ 0.4 \\
 & OLMo-2-1B & 11.7 $\pm$ 1.6 & 20.2 $\pm$ 1.7 & 10.5 $\pm$ 0.3 & 13.4 $\pm$ 2.3 & 22.2 $\pm$ 2.1 & 10.7 $\pm$ 0.7 & 15.2 $\pm$ 0.6 & 13.3 $\pm$ 2.1 & 22.2 $\pm$ 0.4 \\
\midrule
Mean acc. (\%) &  & 13.9 & 21.7 & 16.0 & 16.1 & 19.1 & 13.9 & 25.5 & 24.3 & 26.9 \\
\midrule
Mean tokens &  & 4,820 & 1,403 & 1,076 & 4,616 & 2,585 & 4,570 & 1,429 & 1,218 & 813 \\
\midrule
Mean markers &  & 22.0 & 1.7 & 15.0 & 18.2 & 5.2 & 16.2 & 3.5 & 1.6 & 0.4 \\
\bottomrule
\end{tabular}}

\caption{MATH-500 accuracy (\%) of each method for turning off thinking in the $\lambda=0$ adapter; the bottom three rows give the nine-cell mean accuracy, token count, and frequency of reflective markers. Seed region removed removes the two-seed region; Other removed (seed region) removes the same energy from other directions outside that region; Top removed ($\tau$ = 0.6) removes the top singular components of the update by the criterion of Gold-Switch ($\tau=0.6$). Length-matched cut cuts each cell at the mean length of its adapter with the shared region removed (1,076 tokens averaged over the nine cells) and has the model write the answer. Merged with $\lambda$ = 1 and $\lambda$ = 1 require a separate adapter trained on short-solution data, and Untrained is the untrained student with three different generation seeds. Among the methods that turn thinking off without additional training, removing the two-seed region leaves the fewest reflective markers and is the most accurate, but it is similar to or below the untrained student. Values after $\pm$ are standard deviations across three seeds.}
\label{tab:14}
\end{table*}
